\documentclass[a4paper,14pt]{extarticle}

\usepackage[T2A]{fontenc}
\usepackage[utf8]{inputenc}
\usepackage[russian]{babel}
\usepackage{tempora}
\usepackage[left=3cm, right=1cm, top=2cm, bottom=2cm, footskip=1.25cm]{geometry}
\usepackage{setspace}
\onehalfspacing
\usepackage{indentfirst}
\setlength{\parindent}{1.25cm}
\usepackage{hyphenat}
\sloppy

\usepackage{fancyhdr}
\pagestyle{fancy}
\fancyhf{}
\fancyfoot[C]{\thepage}
\renewcommand{\headrulewidth}{0pt}
\renewcommand{\footrulewidth}{0pt}
\setlength{\headheight}{0pt}

\usepackage{array}
\usepackage{longtable}
\usepackage{tabularx}
\usepackage{booktabs}
\usepackage{graphicx}
\usepackage{float}

\usepackage{listings}
\lstset{
  basicstyle=\small\ttfamily,
  breaklines=true,
  frame=none,
  numbers=none,
  tabsize=4,
  xleftmargin=1.25cm,
  showstringspaces=false,
}

\usepackage{tikz}
\usetikzlibrary{shapes.geometric, arrows.meta, positioning}

\usepackage[labelsep=endash]{caption}
\captionsetup[figure]{justification=centering, font={normalsize}, skip=5pt, position=below}
\captionsetup[table]{justification=raggedleft, singlelinecheck=false, font={normalsize}, skip=5pt, position=above}

\usepackage[hidelinks]{hyperref}

\usepackage{tocloft}
\renewcommand{\cftsecleader}{\cftdotfill{\cftdotsep}}

\usepackage{titlesec}

\newcommand{\structsection}[1]{%
  \clearpage
  \phantomsection
  \begin{center}
    \fontsize{16pt}{20pt}\selectfont\bfseries\MakeUppercase{#1}
  \end{center}
  \addcontentsline{toc}{section}{#1}
  \vspace{1em}
}

\titleformat{\section}[block]
  {\fontsize{16pt}{20pt}\selectfont\bfseries}
  {\thesection}{0.5em}{}
\titleformat{\subsection}[block]
  {\fontsize{14pt}{18pt}\selectfont\bfseries}
  {\thesubsection}{0.5em}{}
\titleformat{\subsubsection}[block]
  {\fontsize{14pt}{18pt}\selectfont\bfseries}
  {\thesubsubsection}{0.5em}{}

\titlespacing*{\section}{1.25cm}{1em}{0.5em}
\titlespacing*{\subsection}{1.25cm}{1em}{0.5em}
\titlespacing*{\subsubsection}{1.25cm}{1em}{0.5em}

\renewcommand{\thesection}{\arabic{section}}
\renewcommand{\thesubsection}{\thesection.\arabic{subsection}}

\usepackage{enumitem}

\begin{document}

% ============================================================
% ТИТУЛЬНЫЙ ЛИСТ
% ============================================================
\thispagestyle{empty}
\begin{center}

\fontsize{14pt}{18pt}\selectfont

МИНИСТЕРСТВО НАУКИ И ВЫСШЕГО ОБРАЗОВАНИЯ\\РОССИЙСКОЙ ФЕДЕРАЦИИ

\vspace{0.5cm}

ФГБОУ ВО <<ДАГЕСТАНСКИЙ ГОСУДАРСТВЕННЫЙ\\ТЕХНИЧЕСКИЙ УНИВЕРСИТЕТ>>

\vspace{0.5cm}

Факультет информатики и вычислительной техники\\
Кафедра информатики

\vspace{2cm}

\fontsize{16pt}{20pt}\selectfont\bfseries
ОТЧЁТ\\
по производственной практике

\vspace{1cm}

\fontsize{14pt}{18pt}\selectfont\normalfont
Тема:\\
<<Разработка RESTful веб-приложения для управления задачами\\
с использованием языка программирования Go\\
и многослойной архитектуры>>

\end{center}

\vspace{2cm}

\begin{flushright}
\fontsize{14pt}{18pt}\selectfont
Выполнил: студент группы {[\_\_\_]}\\
Берсиров~С.\\[1cm]
Руководитель: {[\_\_\_]}
\end{flushright}

\vfill

\begin{center}
\fontsize{14pt}{18pt}\selectfont
Махачкала 2026
\end{center}

\newpage

% ============================================================
% ИНДИВИДУАЛЬНОЕ ЗАДАНИЕ
% ============================================================
\structsection{Индивидуальное задание}

Разработать RESTful веб-приложение TodoApp для управления задачами на языке программирования Go с использованием многослойной архитектуры, СУБД PostgreSQL, контейнеризации Docker и автоматической документации Swagger.

В ходе производственной практики необходимо:
\begin{itemize}[leftmargin=1.25cm, itemsep=0pt]
  \item[---] реализовать серверную часть приложения на языке Go с применением многослойной архитектуры (Transport~--- Service~--- Repository);
  \item[---] реализовать модули управления пользователями (CRUD), задачами (CRUD) и статистики;
  \item[---] обеспечить хранение данных в СУБД PostgreSQL с управлением схемой через миграции;
  \item[---] реализовать REST~API с автоматической Swagger-документацией;
  \item[---] обеспечить контейнеризацию компонентов с помощью Docker~Compose;
  \item[---] провести функциональное тестирование реализованных эндпоинтов.
\end{itemize}

\subsection*{Требования к разрабатываемому программному обеспечению}
\addcontentsline{toc}{subsection}{Требования к разрабатываемому ПО}

\textbf{Функциональные требования:}
\begin{itemize}[leftmargin=1.25cm, itemsep=0pt]
  \item[---] CRUD-операции с пользователями и задачами;
  \item[---] привязка задач к пользователям через поле author\_user\_id;
  \item[---] автоматическая фиксация времени завершения задачи;
  \item[---] статистика с фильтрацией по пользователю и временному диапазону;
  \item[---] пагинация списков (limit, offset);
  \item[---] валидация входных данных;
  \item[---] Swagger-документация по адресу /swagger/;
  \item[---] веб-интерфейс по адресу /.
\end{itemize}

\textbf{Нефункциональные требования:}
\begin{itemize}[leftmargin=1.25cm, itemsep=0pt]
  \item[---] graceful shutdown с таймаутом 30~секунд;
  \item[---] пул соединений с PostgreSQL (pgx);
  \item[---] структурированное логирование (zap);
  \item[---] поддержка CORS;
  \item[---] оптимистичная блокировка;
  \item[---] контейнеризация через Docker~Compose.
\end{itemize}

% ============================================================
% СОДЕРЖАНИЕ
% ============================================================
\clearpage
\phantomsection
\begin{center}
\fontsize{16pt}{20pt}\selectfont\bfseries СОДЕРЖАНИЕ
\end{center}
\vspace{1em}
\renewcommand{\contentsname}{}
\tableofcontents

% ============================================================
% ВВЕДЕНИЕ
% ============================================================
\structsection{Введение}

Производственная практика является важным этапом подготовки специалиста в области информационных технологий, позволяющим закрепить теоретические знания путём решения практических задач разработки программного обеспечения.

Целью производственной практики является разработка RESTful веб-приложения для управления задачами (TodoApp) на языке программирования Go с применением многослойной архитектуры, СУБД PostgreSQL и технологий контейнеризации Docker.

Задачи практики:
\begin{itemize}[leftmargin=1.25cm, itemsep=0pt]
  \item[---] реализовать серверную часть с разделением на слои Transport, Service и Repository;
  \item[---] реализовать модули управления пользователями, задачами и статистикой;
  \item[---] обеспечить хранение данных в PostgreSQL с миграциями;
  \item[---] реализовать REST~API с Swagger-документацией;
  \item[---] провести функциональное тестирование.
\end{itemize}

Практика выполнена на базе проекта TodoApp~--- веб-приложения с открытым исходным кодом, размещённого в репозитории GitHub. Приложение реализует полный цикл CRUD-операций для пользователей и задач, предоставляет REST~API и веб-интерфейс.

% ============================================================
% 1 АЛГОРИТМ РАБОТЫ ПРИЛОЖЕНИЯ
% ============================================================
\clearpage
\section{Алгоритм работы приложения}

\subsection{Общий алгоритм обработки HTTP-запросов}

Обработка каждого HTTP-запроса проходит через цепочку компонентов, организованных в конвейер. Алгоритм представлен на рисунке~\ref{fig:pipeline}.

\begin{figure}[H]
\centering
\begin{tikzpicture}[
  box/.style={rectangle, draw, minimum width=4.5cm, minimum height=0.7cm, text centered, font=\small},
  arr/.style={-{Stealth[length=2.5mm]}, thick}
]
\node[box] (req) {HTTP-запрос от клиента};
\node[box, below=0.5cm of req] (cors) {Middleware: CORS};
\node[box, below=0.4cm of cors] (rid) {Middleware: RequestID};
\node[box, below=0.4cm of rid] (log) {Middleware: Logger};
\node[box, below=0.4cm of log] (trace) {Middleware: Trace};
\node[box, below=0.4cm of trace] (panic) {Middleware: Panic};
\node[box, below=0.4cm of panic] (router) {Router (ServeMux)};
\node[box, below=0.4cm of router] (handler) {Handler};
\node[box, below=0.4cm of handler] (service) {Service};
\node[box, below=0.4cm of service] (repo) {Repository};
\node[box, below=0.4cm of repo] (db) {PostgreSQL};

\draw[arr] (req) -- (cors);
\draw[arr] (cors) -- (rid);
\draw[arr] (rid) -- (log);
\draw[arr] (log) -- (trace);
\draw[arr] (trace) -- (panic);
\draw[arr] (panic) -- (router);
\draw[arr] (router) -- (handler);
\draw[arr] (handler) -- (service);
\draw[arr] (service) -- (repo);
\draw[arr] (repo) -- (db);
\end{tikzpicture}
\caption{Блок-схема обработки HTTP-запроса}\label{fig:pipeline}
\end{figure}

Middleware применяются ко всем маршрутам сервера и выполняются в порядке регистрации. Каждый middleware оборачивает следующий обработчик, формируя цепочку вызовов.

\subsection{Алгоритм создания задачи}

Алгоритм создания задачи включает следующие шаги:
\begin{enumerate}[leftmargin=1.25cm, itemsep=0pt]
  \item клиент отправляет POST-запрос на /api/v1/tasks с телом в формате JSON;
  \item транспортный слой десериализует тело и выполняет валидацию через go-playground/validator;
  \item сервис вызывает конструктор domain.CreateTask, генерирующий UUID, version~=~1, completed~=~false, created\_at~=~now();
  \item вызывается метод Validate() для проверки инвариантов;
  \item репозиторий выполняет INSERT INTO tasks \ldots{} RETURNING;
  \item транспортный слой формирует JSON-ответ с HTTP~201.
\end{enumerate}

\subsection{Алгоритм обновления задачи (PATCH с оптимистичной блокировкой)}

Обновление задачи реализовано через PATCH-запросы:
\begin{enumerate}[leftmargin=1.25cm, itemsep=0pt]
  \item клиент отправляет PATCH с частичными данными (тип Nullable[T]);
  \item сервис получает текущую задачу из БД;
  \item применяется ApplyPatch (работа с копией объекта);
  \item при completed~=~true автоматически устанавливается completed\_at;
  \item UPDATE \ldots{} WHERE id~=~\$1 AND version~=~\$2;
  \item при конфликте версии~--- HTTP~409.
\end{enumerate}

% ============================================================
% 2 МОДУЛЬНАЯ СТРУКТУРА
% ============================================================
\clearpage
\section{Модульная структура приложения}

Приложение организовано по принципу feature-based structure. Модульная диаграмма представлена на рисунке~\ref{fig:modules}.

\begin{figure}[H]
\centering
\begin{tikzpicture}[
  box/.style={rectangle, draw, minimum width=2.5cm, minimum height=0.8cm, text centered, font=\small},
  arr/.style={-{Stealth[length=2mm]}, thick}
]
\node[box] (main) {main.go};
\node[box, below left=1cm and 0.5cm of main] (users) {users};
\node[box, below=1cm of main] (tasks) {tasks};
\node[box, below right=1cm and 0.5cm of main] (stats) {statistics};
\node[box, right=0.3cm of stats] (web) {web};
\node[box, below=1.5cm of tasks] (core) {core};

\draw[arr] (main) -- (users);
\draw[arr] (main) -- (tasks);
\draw[arr] (main) -- (stats);
\draw[arr] (main) -- (web);
\draw[arr] (users) -- (core);
\draw[arr] (tasks) -- (core);
\draw[arr] (stats) -- (core);
\draw[arr] (web) -- (core);
\end{tikzpicture}
\caption{Модульная диаграмма приложения}\label{fig:modules}
\end{figure}

\subsection{Модуль core (ядро)}

Содержит общие компоненты: domain (сущности Task, User, Nullable[T]), errors, config, logger, пул соединений PostgreSQL, HTTP-сервер, middleware (CORS, RequestID, Logger, Trace, Panic), обработку запросов и ответов.

\subsection{Модуль users (пользователи)}

Реализует CRUD-операции: CreateUser, GetUsers, GetUser, PatchUser, DeleteUser. Интерфейс UsersRepository определён в пакете сервиса по принципу Dependency Inversion.

\subsection{Модуль tasks (задачи)}

Реализует CRUD-операции: CreateTask, GetTasks, GetTask, PatchTask, DeleteTask. Поддерживает фильтрацию по user\_id и пагинацию.

\subsection{Модуль statistics (статистика)}

Реализует GetStatistics с параметрами фильтрации. Интерфейс StatisticsRepository содержит только GetTasks~--- пример принципа Interface Segregation.

\subsection{Модуль web (веб-интерфейс)}

Обслуживает HTML-файл веб-интерфейса через обработчик GetMainPage.

% ============================================================
% 3 РЕАЛИЗАЦИЯ
% ============================================================
\clearpage
\section{Реализация}

\subsection{Доменный слой}

Доменный слой содержит сущности и бизнес-правила. Структура сущности Task:

\begin{lstlisting}[language=Go]
type Task struct {
    ID           uuid.UUID
    Version      int
    Title        string
    Description  *string
    Completed    bool
    CreatedAt    time.Time
    CompletedAt  *time.Time
    AuthorUserID uuid.UUID
}
\end{lstlisting}

Конструктор CreateTask генерирует UUID и начальные значения:

\begin{lstlisting}[language=Go]
func CreateTask(
    title string,
    description *string,
    authorUserID uuid.UUID,
) Task {
    var (
        id          = uuid.New()
        version     = 1
        completed   = false
        createdAt   = time.Now()
        completedAt *time.Time = nil
    )
    return NewTask(id, version, title, description,
        completed, createdAt, completedAt, authorUserID)
}
\end{lstlisting}

Метод Validate() проверяет инварианты с использованием len([]rune(\ldots)):

\begin{lstlisting}[language=Go]
func (t *Task) Validate() error {
    titleLen := len([]rune(t.Title))
    if titleLen < 1 || titleLen > 100 {
        return fmt.Errorf(
            "invalid Title len: %d: %w",
            titleLen, core_errors.ErrInvalidArgument)
    }
    // ... проверка Description, Completed/CompletedAt
    return nil
}
\end{lstlisting}

Метод ApplyPatch работает с копией для транзакционности:

\begin{lstlisting}[language=Go]
func (t *Task) ApplyPatch(patch TaskPatch) error {
    if err := patch.Validate(); err != nil {
        return err
    }
    tmp := *t
    if patch.Title.Set {
        tmp.Title = *patch.Title.Value
    }
    if patch.Completed.Set {
        tmp.Completed = *patch.Completed.Value
        if tmp.Completed {
            completedAt := time.Now()
            tmp.CompletedAt = &completedAt
        } else {
            tmp.CompletedAt = nil
        }
    }
    if err := tmp.Validate(); err != nil {
        return err
    }
    *t = tmp
    return nil
}
\end{lstlisting}

Сущность User:

\begin{lstlisting}[language=Go]
type User struct {
    ID          uuid.UUID
    Version     int
    FullName    string
    PhoneNumber *string
}
\end{lstlisting}

Валидация телефона (формат +цифры):

\begin{lstlisting}[language=Go]
func (u *User) Validate() error {
    fullNameLen := len([]rune(u.FullName))
    if fullNameLen < 3 || fullNameLen > 100 {
        return fmt.Errorf("invalid FullName len: %d: %w",
            fullNameLen, core_errors.ErrInvalidArgument)
    }
    if u.PhoneNumber != nil {
        re := regexp.MustCompile(`^\+[0-9]+$`)
        if !re.MatchString(*u.PhoneNumber) {
            return fmt.Errorf("invalid PhoneNumber: %w",
                core_errors.ErrInvalidArgument)
        }
    }
    return nil
}
\end{lstlisting}

\subsection{Сервисный слой}

Интерфейс репозитория определён в пакете сервиса:

\begin{lstlisting}[language=Go]
type TasksRepository interface {
    SaveTask(ctx context.Context,
        task domain.Task) (domain.Task, error)
    GetTasks(ctx context.Context, userID *uuid.UUID,
        limit *int, offset *int) ([]domain.Task, error)
    GetTask(ctx context.Context,
        id uuid.UUID) (domain.Task, error)
    DeleteTask(ctx context.Context,
        id uuid.UUID) error
    UpdateTask(ctx context.Context,
        task domain.Task) (domain.Task, error)
}
\end{lstlisting}

Создание задачи в сервисе:

\begin{lstlisting}[language=Go]
func (s *TasksService) CreateTask(
    ctx context.Context,
    title string,
    description *string,
    authorUserID uuid.UUID,
) (domain.Task, error) {
    task := domain.CreateTask(title, description,
        authorUserID)
    if err := task.Validate(); err != nil {
        return domain.Task{}, err
    }
    return s.tasksRepository.SaveTask(ctx, task)
}
\end{lstlisting}

\subsection{Слой репозитория}

Сохранение задачи в PostgreSQL:

\begin{lstlisting}[language=Go]
func (r *TasksRepository) SaveTask(
    ctx context.Context, task domain.Task,
) (domain.Task, error) {
    row := r.pool.QueryRow(ctx,
        `INSERT INTO tasks
            (id, version, title, description,
             completed, created_at, completed_at,
             author_user_id)
         VALUES ($1,$2,$3,$4,$5,$6,$7,$8)
         RETURNING id, version, title, description,
             completed, created_at, completed_at,
             author_user_id`,
        task.ID, task.Version, task.Title,
        task.Description, task.Completed,
        task.CreatedAt, task.CompletedAt,
        task.AuthorUserID)
    var model TaskModel
    err := row.Scan(&model.ID, &model.Version,
        &model.Title, &model.Description,
        &model.Completed, &model.CreatedAt,
        &model.CompletedAt, &model.AuthorUserID)
    if err != nil {
        return domain.Task{}, err
    }
    return model.ToDomain(), nil
}
\end{lstlisting}

Обновление с оптимистичной блокировкой:

\begin{lstlisting}[language=Go]
`UPDATE tasks
 SET version = version + 1,
     title = $3, description = $4,
     completed = $5, completed_at = $6
 WHERE id = $1 AND version = $2
 RETURNING ...`
\end{lstlisting}

\subsection{Транспортный слой (HTTP)}

Обработчик создания пользователя:

\begin{lstlisting}[language=Go]
type CreateUserRequest struct {
    FullName    string  `json:"full_name"
        validate:"required,min=3,max=100"`
    PhoneNumber *string `json:"phone_number"
        validate:"omitempty,min=10,max=15,
        startswith=+"`
}

func (h *UsersHTTPHandler) CreateUser(
    rw http.ResponseWriter, r *http.Request,
) {
    ctx := r.Context()
    log := core_logger.FromContext(ctx)
    handler := core_http_response.
        NewHTTPResponseHandler(log, rw)

    var request CreateUserRequest
    if err := core_http_request.
        DecodeAndValidateRequest(r, &request);
        err != nil {
        handler.ErrorResponse(err, "validation error")
        return
    }
    userDomain, err := h.usersService.CreateUser(
        ctx, request.FullName, request.PhoneNumber)
    if err != nil {
        handler.ErrorResponse(err, "create error")
        return
    }
    response := CreateUserResponse(
        userDTOFromDomain(userDomain))
    handler.JSONResponse(response, http.StatusCreated)
}
\end{lstlisting}

\subsection{Конфигурация и запуск}

Точка входа~--- main.go. Ручное внедрение зависимостей:

\begin{lstlisting}[language=Go]
func main() {
    cfg := core_config.NewConfigMust()
    ctx, cancel := signal.NotifyContext(
        context.Background(),
        syscall.SIGINT, syscall.SIGTERM)
    defer cancel()

    logger, _ := core_logger.NewLogger(
        core_logger.NewConfigMust())
    defer logger.Close()

    postgresPool, _ := core_pgx_pool.NewPool(
        ctx, core_pgx_pool.NewConfigMust())
    defer postgresPool.Close()

    // Repository -> Service -> Handler
    usersRepo := users_postgres_repository.
        NewUsersRepository(postgresPool)
    usersSvc := users_service.
        NewUsersService(usersRepo)
    usersHTTP := users_transport_http.
        NewUsersHTTPHandler(usersSvc)
    // аналогично tasks, statistics, web...

    httpServer := core_http_server.NewHTTPServer(
        httpConfig, logger,
        core_http_middleware.CORS(...),
        core_http_middleware.RequestID(),
        core_http_middleware.Logger(logger),
        core_http_middleware.Trace(),
        core_http_middleware.Panic())

    router := core_http_server.NewAPIVersionRouter(
        core_http_server.ApiVersion1)
    router.AddRoutes(usersHTTP.Routes()...)
    router.AddRoutes(tasksHTTP.Routes()...)
    router.AddRoutes(statsHTTP.Routes()...)

    httpServer.RegisterAPIRouters(router)
    httpServer.RegisterSwagger()
    httpServer.Run(ctx)
}
\end{lstlisting}

% ============================================================
% 4 ТЕСТИРОВАНИЕ
% ============================================================
\clearpage
\section{Тестирование}

Функциональное тестирование проведено через Swagger~UI и веб-интерфейс. Результаты представлены в таблице~\ref{tab:tests}.

\begin{table}[H]
\caption{Результаты функционального тестирования}\label{tab:tests}
\centering
\small
\begin{tabularx}{\textwidth}{|c|X|c|c|}
\hline
\textbf{№} & \textbf{Тестовый сценарий} & \textbf{Ожидание} & \textbf{Статус} \\
\hline
1 & Создание пользователя с валидными данными & HTTP 201 & Пройден \\
\hline
2 & Создание пользователя с невалидным телефоном & HTTP 400 & Пройден \\
\hline
3 & Получение списка пользователей & HTTP 200 & Пройден \\
\hline
4 & Создание задачи с существующим автором & HTTP 201 & Пройден \\
\hline
5 & Создание задачи без автора & HTTP 400 & Пройден \\
\hline
6 & Обновление статуса задачи (completed = true) & HTTP 200 & Пройден \\
\hline
7 & Удаление задачи & HTTP 204 & Пройден \\
\hline
8 & Получение несуществующей задачи & HTTP 404 & Пройден \\
\hline
9 & Получение статистики & HTTP 200 & Пройден \\
\hline
10 & Обновление с конфликтом версии & HTTP 409 & Пройден \\
\hline
\end{tabularx}
\end{table}

Рисунок~\ref{fig:swagger}~--- скриншот Swagger~UI (вставить скриншот).

\begin{figure}[H]
\centering
\fbox{\parbox{10cm}{\centering\vspace{3cm}\textit{Вставить скриншот Swagger UI\\http://127.0.0.1:5050/swagger/}\vspace{3cm}}}
\caption{Скриншот Swagger~UI}\label{fig:swagger}
\end{figure}

Рисунок~\ref{fig:webui}~--- скриншот веб-интерфейса (вставить скриншот).

\begin{figure}[H]
\centering
\fbox{\parbox{10cm}{\centering\vspace{3cm}\textit{Вставить скриншот веб-интерфейса TodoApp\\http://127.0.0.1:5050/}\vspace{3cm}}}
\caption{Скриншот веб-интерфейса TodoApp}\label{fig:webui}
\end{figure}

Все десять тестовых сценариев пройдены успешно.

% ============================================================
% 5 ОЦЕНКА РЕЗУЛЬТАТОВ
% ============================================================
\section{Оценка результатов}

Реализованные функциональные возможности:
\begin{itemize}[leftmargin=1.25cm, itemsep=0pt]
  \item[---] полный CRUD для пользователей и задач с валидацией;
  \item[---] автоматическое управление временем завершения задачи;
  \item[---] оптимистичная блокировка;
  \item[---] статистика с фильтрацией;
  \item[---] пагинация списков;
  \item[---] Swagger-документация;
  \item[---] веб-интерфейс.
\end{itemize}

Реализованные нефункциональные требования:
\begin{itemize}[leftmargin=1.25cm, itemsep=0pt]
  \item[---] graceful shutdown (30~с);
  \item[---] пул соединений pgx;
  \item[---] логирование zap;
  \item[---] CORS;
  \item[---] контейнеризация Docker~Compose.
\end{itemize}

Все требования технического задания реализованы в полном объёме.

% ============================================================
% ЗАКЛЮЧЕНИЕ
% ============================================================
\structsection{Заключение}

В ходе производственной практики разработано RESTful веб-приложение TodoApp на языке Go с многослойной архитектурой.

Выполнены следующие работы:
\begin{itemize}[leftmargin=1.25cm, itemsep=0pt]
  \item[---] реализована серверная часть с разделением на слои Transport, Service, Repository;
  \item[---] реализованы модули: users, tasks, statistics, web;
  \item[---] обеспечено хранение данных в PostgreSQL с миграциями;
  \item[---] реализован REST~API с 11~эндпоинтами и Swagger-документацией;
  \item[---] применены паттерны: DI, Repository, Optimistic Locking;
  \item[---] контейнеризация через Docker~Compose;
  \item[---] проведено тестирование: 10~сценариев пройдены успешно.
\end{itemize}

Полученные навыки: разработка на Go, проектирование REST~API, работа с PostgreSQL через pgx, применение многослойной архитектуры и SOLID, контейнеризация Docker, Git.

% ============================================================
% СПИСОК ИСПОЛЬЗОВАННЫХ ИСТОЧНИКОВ
% ============================================================
\structsection{Список использованных источников}

\begin{enumerate}[leftmargin=1.25cm, itemsep=2pt, label=\arabic*.]
  \item Донован,~А. Язык программирования Go~/~А.~Донован, Б.~Керниган.~--- М.~: Вильямс, 2021.~--- 432~с.
  \item Мартин,~Р. Чистая архитектура~/~Р.~Мартин.~--- СПб.~: Питер, 2018.~--- 352~с.
  \item Fielding,~R.~T. Architectural Styles and the Design of Network-based Software Architectures~: дис.~/~R.~T.~Fielding.~--- University of California, Irvine, 2000.~--- 162~p.
  \item PostgreSQL~18 Documentation~[Электронный ресурс].~--- Режим доступа: https://www.postgresql.org/docs/18/ (дата обращения: 01.06.2026).
  \item The Go Programming Language Specification~[Электронный ресурс].~--- Режим доступа: https://go.dev/ref/spec (дата обращения: 01.06.2026).
  \item Docker Documentation~[Электронный ресурс].~--- Режим доступа: https://docs.docker.com/ (дата обращения: 01.06.2026).
  \item OpenAPI Specification~v3.0~[Электронный ресурс].~--- Режим доступа: https://swagger.io/specification/ (дата обращения: 01.06.2026).
  \item ГОСТ~7.32-2017. Отчёт о научно-исследовательской работе.~--- М.~: Стандартинформ, 2017.
  \item ГОСТ~Р~7.0.5-2008. Библиографическая ссылка.~--- М.~: Стандартинформ, 2008.
  \item Gamma,~E. Design Patterns~/~E.~Gamma, R.~Helm, R.~Johnson, J.~Vlissides.~--- Addison-Wesley, 1994.~--- 395~p.
  \item Фаулер,~М. Шаблоны корпоративных приложений~/~М.~Фаулер.~--- М.~: Вильямс, 2020.~--- 544~с.
  \item Richardson,~L. RESTful Web Services~/~L.~Richardson, S.~Ruby.~--- O'Reilly Media, 2007.~--- 454~p.
  \item pgx~--- PostgreSQL Driver for Go~[Электронный ресурс].~--- Режим доступа: https://github.com/jackc/pgx (дата обращения: 01.06.2026).
  \item Zap~--- Structured logging in Go~[Электронный ресурс].~--- Режим доступа: https://github.com/uber-go/zap (дата обращения: 01.06.2026).
  \item golang-migrate~[Электронный ресурс].~--- Режим доступа: https://github.com/golang-migrate/migrate (дата обращения: 01.06.2026).
\end{enumerate}

% ============================================================
% ПРИЛОЖЕНИЯ
% ============================================================
\structsection{Приложение А}
\begin{center}(обязательное)\end{center}
\begin{center}\textbf{Исходный код main.go}\end{center}

\begin{lstlisting}[language=Go]
package main

import (
    "context"
    "fmt"
    "os"
    "os/signal"
    "syscall"
    "time"

    core_config "github.com/milas1221/golang-todoapp/internal/core/config"
    core_logger "github.com/milas1221/golang-todoapp/internal/core/logger"
    core_pgx_pool "github.com/milas1221/golang-todoapp/internal/core/repository/postgres/pool/pgx"
    core_http_middleware "github.com/milas1221/golang-todoapp/internal/core/transport/http/middleware"
    core_http_server "github.com/milas1221/golang-todoapp/internal/core/transport/http/server"
    statistics_postgres_repository "github.com/milas1221/golang-todoapp/internal/features/statistics/repository/postgres"
    statistics_service "github.com/milas1221/golang-todoapp/internal/features/statistics/service"
    statistics_transport_http "github.com/milas1221/golang-todoapp/internal/features/statistics/transport/http"
    tasks_postgres_repository "github.com/milas1221/golang-todoapp/internal/features/tasks/repository/postgres"
    tasks_service "github.com/milas1221/golang-todoapp/internal/features/tasks/service"
    tasks_transport_http "github.com/milas1221/golang-todoapp/internal/features/tasks/transport/http"
    users_postgres_repository "github.com/milas1221/golang-todoapp/internal/features/users/repository/postgres"
    users_service "github.com/milas1221/golang-todoapp/internal/features/users/service"
    users_transport_http "github.com/milas1221/golang-todoapp/internal/features/users/transport/http"
    web_fs_repository "github.com/milas1221/golang-todoapp/internal/features/web/repository/file_system"
    web_service "github.com/milas1221/golang-todoapp/internal/features/web/service"
    web_transport_http "github.com/milas1221/golang-todoapp/internal/features/web/transport/http"
    "go.uber.org/zap"
    _ "github.com/milas1221/golang-todoapp/docs"
)

func main() {
    cfg := core_config.NewConfigMust()
    time.Local = cfg.TimeZone
    ctx, cancel := signal.NotifyContext(
        context.Background(),
        syscall.SIGINT, syscall.SIGTERM)
    defer cancel()

    logger, err := core_logger.NewLogger(
        core_logger.NewConfigMust())
    if err != nil {
        fmt.Println("failed to init logger:", err)
        os.Exit(1)
    }
    defer logger.Close()

    postgresPool, err := core_pgx_pool.NewPool(
        ctx, core_pgx_pool.NewConfigMust())
    if err != nil {
        logger.Fatal("failed to init pool",
            zap.Error(err))
    }
    defer postgresPool.Close()

    usersRepository := users_postgres_repository.
        NewUsersRepository(postgresPool)
    usersService := users_service.
        NewUsersService(usersRepository)
    usersTransportHTTP := users_transport_http.
        NewUsersHTTPHandler(usersService)

    tasksRepository := tasks_postgres_repository.
        NewTasksRepository(postgresPool)
    tasksService := tasks_service.
        NewTasksService(tasksRepository)
    tasksTransportHTTP := tasks_transport_http.
        NewTasksHTTPHandler(tasksService)

    statisticsRepository := statistics_postgres_repository.
        NewStatisticsRepository(postgresPool)
    statisticsService := statistics_service.
        NewStatisticsService(statisticsRepository)
    statisticsTransportHTTP := statistics_transport_http.
        NewStatisticsHTTPHandler(statisticsService)

    webRepository := web_fs_repository.
        NewWebRepository()
    webService := web_service.
        NewWebService(webRepository)
    webTransportHTTP := web_transport_http.
        NewWebHTTPHandler(webService)

    httpConfig := core_http_server.NewConfigMust()
    httpServer := core_http_server.NewHTTPServer(
        httpConfig, logger,
        core_http_middleware.CORS(
            httpConfig.AllowedOrigins),
        core_http_middleware.RequestID(),
        core_http_middleware.Logger(logger),
        core_http_middleware.Trace(),
        core_http_middleware.Panic())

    apiV1 := core_http_server.NewAPIVersionRouter(
        core_http_server.ApiVersion1)
    apiV1.AddRoutes(usersTransportHTTP.Routes()...)
    apiV1.AddRoutes(tasksTransportHTTP.Routes()...)
    apiV1.AddRoutes(
        statisticsTransportHTTP.Routes()...)

    httpServer.RegisterAPIRouters(apiV1)
    httpServer.RegisterRoutes(
        webTransportHTTP.Routes()...)
    httpServer.RegisterSwagger()

    if err := httpServer.Run(ctx); err != nil {
        logger.Error("server error",
            zap.Error(err))
    }
}
\end{lstlisting}

\structsection{Приложение Б}
\begin{center}(обязательное)\end{center}
\begin{center}\textbf{Исходный код доменных моделей}\end{center}

Файл internal/core/domain/task.go и internal/core/domain/user.go приведены в разделе~3.1 настоящего отчёта.

\structsection{Приложение В}
\begin{center}(обязательное)\end{center}
\begin{center}\textbf{Файл docker-compose.yaml}\end{center}

\begin{lstlisting}
services:
  todoapp:
    build:
      context: ${PROJECT_ROOT}
      dockerfile: ${PROJECT_ROOT}/cmd/todoapp/Dockerfile
    container_name: todoapp
    env_file:
      - ${PROJECT_ROOT}/.env
    environment:
      - HTTP_ADDR=:5050
      - POSTGRES_HOST=todoapp-postgres
      - LOGGER_FOLDER=/app/out/logs
    ports:
      - "5050:5050"
    restart: always

  todoapp-postgres:
    image: postgres:18.1-bookworm
    container_name: todoapp-env-postgres
    environment:
      - POSTGRES_USER=${POSTGRES_USER}
      - POSTGRES_PASSWORD=${POSTGRES_PASSWORD}
      - POSTGRES_DB=${POSTGRES_DB}
    volumes:
      - ${PROJECT_ROOT}/out/pgdata:/var/lib/postgresql
    restart: always

  todoapp-postgres-migrate:
    image: migrate/migrate:v4.19.1
    volumes:
      - ${PROJECT_ROOT}/migrations:/migrations

  port-forwarder:
    image: alpine/socat:1.8.0.3
    container_name: todoapp-env-port-forwarder
    command: >-
      tcp-listen:5433,fork,reuseaddr
      tcp-connect:todoapp-postgres:5432
    ports:
      - "127.0.0.1:5433:5433"

  swagger:
    build:
      context: ${PROJECT_ROOT}
      dockerfile_inline: |
        FROM golang:1.24.1-alpine
        RUN apk add --no-cache git
        RUN go install github.com/swaggo/swag/cmd/swag@v1.16.6
        WORKDIR /code
        ENTRYPOINT ["swag"]
    volumes:
      - ${PROJECT_ROOT}:/code
\end{lstlisting}

\end{document}
